% =====================================================================
% SHIFT-SURV — revisione 2026-08-19
% Blocchi sostitutivi per: Abstract, Introduction, coda Results real-data,
% Ablation/limitations, Discussion.
%
% Ogni blocco e' delimitato da ==== REPLACE ... ==== END.
% Sostituire l'intero blocco corrispondente nel main .tex su Overleaf.
%
% NUMERI: tutti verificati contro outputs/tables/*.csv il 2026-08-19,
% TRANNE quelli marcati [VERIFICARE] nel blocco Cox, il cui CSV
% (i3lung_survival_transfers.csv) non e' nel repo perche' i dati I3LUNG
% sono privati. Rilanciare:
%   python -m experiments.run_survival_real --data_dir <I3LUNG_DIR>
% =====================================================================


% ==== REPLACE: abstract ==============================================
\begin{abstract}
Risk models for time-to-event outcomes are increasingly trained on pooled
multi-center cohorts and deployed at hospitals that did not contribute to
training. There, differences in patient case mix, measurement protocols,
treatment pathways, and outcome recording degrade their performance.
Unsupervised domain adaptation and domain generalization are the standard
remedies, but on real cohorts a poor target result cannot be attributed: it is
not possible to tell whether the method underperformed or whether the transfer
was recoverable at all from the information available.

We introduce SHIFT-SURV, a controlled benchmark for fixed-horizon survival-risk
prediction. A latent-factor Weibull proportional-hazards model generates centers
under six independently varied shift mechanisms---covariate, measurement,
spurious-association, baseline-hazard, label-prior, and concept---and evaluates
methods against two references: a latent oracle fitted on the true invariant
factors, and an observed-data oracle fitted on the target covariates with target
labels. The gap between them separates information destroyed by the observation
process from performance a method leaves unrecovered.

No adaptation or generalization method improves discrimination over empirical
risk minimization (ERM) in any regime, including the regimes where the
observed-data oracle shows the signal is still present in the target covariates.
Under measurement shift the best method remains $0.10$ AUROC below that
reference, and the gap does not close as the per-center sample grows from $250$
to $2000$ patients, indicating an identifiability limit rather than a power
limit. Calibration responds to the same shifts along an independent axis and is
recoverable without target labels only under label-prior shift. Both patterns
reappear across $30$ transfers in the six-center I3LUNG lung-cancer cohort and in
the public Rotterdam/GBSG2 registry, and a Cox analysis on the full censored
follow-up reproduces the degradation, ruling out the fixed-horizon
dichotomization as its cause.
\end{abstract}
% ==== END abstract ===================================================


% ==== REPLACE: \section{Introduction} (intero) ========================
\section{Introduction}

Time-to-event risk prediction is widely used in oncology to estimate clinically
relevant outcomes such as disease progression or
death~\cite{delgado2021clinicalendpoints}. These models are increasingly
developed on data pooled across hospitals, registries, and multi-center
consortia, and are then transported to centers that did not contribute to
training. At the new center, differences in patient populations, measurement
protocols, treatment pathways, and outcome recording degrade predictive
performance. This failure to transport is a central obstacle to the clinical
use of such models across sites.

Several families of methods address this problem. Survival transfer methods use
labeled data from one or more source cohorts and may additionally exploit
unlabeled covariates from the target
center~\cite{shaker2023multisource,Bel19,li2023transfer}. Domain generalization
follows a different protocol: no target data are available during training, and
variation across the observed source environments is used to learn a predictor
intended to transfer to an unseen center~\cite{pfisterer2022evaluating,
fan2024stablecox}. Recalibration methods address a related but distinct problem,
updating predicted risks after transport~\cite{ye2025multicalibration}. Yet
empirical studies repeatedly find that these specialized methods do not
consistently outperform empirical risk minimization (ERM) under realistic
shift~\cite{gulrajani2021search,koh2021wilds,guo2022evaluation,zhang2021empirical}.

Such findings are hard to interpret, because real clinical data do not identify
the mechanism responsible for the failure. Differences between centers may
reflect changes in case mix, in the measurement process, in center-specific
associations that do not persist after transport, in the baseline event rate, in
the marginal distribution of the endpoint, or in the conditional relationship
between predictors and outcome. These mechanisms have different implications for
transportability---some leave a target rule that is still recoverable from the
observed covariates, others do not---but in observational cohorts they are
entangled and unobserved. A method that fails to improve on ERM may therefore be
inadequate, or it may be facing a transfer that no label-free method could
solve. Nothing in the data distinguishes the two.

We introduce SHIFT-SURV, a controlled benchmark for fixed-horizon survival-risk
prediction that makes this distinction measurable. Tabular biomarker-like data
are generated from a latent-factor Weibull proportional-hazards model, and six
representative clinical shift mechanisms are varied one at a time under known
ground truth, following simulation frameworks for evaluating robustness under
specified shift~\cite{lasko2024why,zhang2021empirical,coombes2021umpire,
sun2021recovering} and the established taxonomy of dataset
shift~\cite{morenotorres2012unifying}. Performance is read against two
references: a \emph{latent oracle} fitted on the true invariant factors, and an
\emph{observed-data oracle} fitted on the target covariates using target labels.
The first states how much outcome-relevant signal exists before the observation
process; the second states how much of it survives into the recorded covariates.
Their gap separates information that observation has destroyed---which no
method can recover---from performance that a practical method leaves on the
table. The generator produces continuous event times, and the fixed-horizon
binary endpoint $Y^{(\tau)}=\mathbf{1}[T\leq\tau]$ is derived from them; the
event times are retained so that the transport conclusions can be re-checked on
the full censored follow-up (Section~\ref{sec:real}).

The benchmark was designed with clinical input from the oncologists of the
I3LUNG consortium, a six-center lung-cancer cohort. The six-month horizon is the
endpoint they use to identify short survivors, and the eleven clinical variables
of that cohort are the ones they identified as routinely available at that
decision point. The shift mechanisms were chosen to reflect differences in
patient populations, clinical pathways, and measurement protocols observed
across the participating centers, and the two-reference decomposition was
motivated by inspecting where risk models trained on I3LUNG failed, rather than
chosen for methodological convenience.

Across the six regimes, no adaptation or generalization method recovers
discrimination beyond ERM---not even where the observed-data oracle shows the
signal remains available---while calibration degrades along an independent axis
and is correctable without labels in exactly one regime. The same two patterns
hold on both real cohorts. Our contributions are:

\begin{itemize}
    \item \textbf{SHIFT-SURV}, a controlled benchmark for fixed-horizon
    survival-risk transport, in which six representative clinical shift
    mechanisms are varied independently under known ground truth;

    \item a \textbf{two-reference decomposition} separating information
    unavailable in the observed covariates from the margin a method leaves
    unrecovered, making the difference between method failure and unrecoverable
    transfer measurable rather than assumed;

    \item an \textbf{evaluation of representative adaptation,
    domain-generalization, and recalibration methods} across the controlled
    regimes and on two real multi-center cohorts, with the real-cohort
    conclusion verified on the full censored follow-up.
\end{itemize}
% ==== END Introduction ================================================


% ==== INSERT: nuovo paragrafo in coda a §I3LUNG ======================
% Va DOPO il paragrafo "Calibration showed a different pattern..." e PRIMA
% di \subsubsection{Public breast-cancer registry}.
% [VERIFICARE] i cinque numeri contro l'output di run_survival_real.py.

\textbf{Survival validity check.}
The analysis above uses a binary endpoint at $\tau=6$ months, which discards
the $103$ patients whose status at that horizon is unknown and ignores the
information in the censored follow-up. To confirm that the observed degradation
is not produced by this dichotomization, we refitted all $30$ ordered center
pairs as Cox proportional-hazards models on the full follow-up, using overall
survival time and the death indicator. This retains all $2074$ patients and
accommodates the $\approx\!30\%$ censoring. Discrimination degrades from source
to target in $23$ of the $30$ transfers, with a median C-index gap of $+0.034$
(range $-0.083$ to $+0.105$; Wilcoxon $p=0.004$). The size of the gap is
uncorrelated with the fraction of patients each center loses to the fixed
horizon (Spearman $\rho=-0.32$, $p=0.09$), and the sign of that association is
opposite to what differential exclusion would produce, with the largest gaps
occurring at $0\%$ and $12.8\%$ exclusion. The transport failure is therefore a
property of the data and not of the endpoint definition. This check compares
only ERM-equivalent models, since the adversarial baselines have no
survival-loss counterpart.
% ==== END survival check =============================================


% ==== REPLACE: \section{Ablation studies and limitations} (intero) ====
\section{Ablation studies and limitations}
\label{sec:ablation}

\textbf{Sample size.}
The discrimination result is not a power limitation. Sweeping the per-center
sample size from $250$ to $2000$ patients (Fig.~\ref{fig:nsweep}a), the best
adaptation method never improves on ERM by more than $0.007$ AUROC at any size.
The small advantage it holds at $250$ patients under covariate shift shrinks as
patients are added rather than growing, and under spurious shift it falls
further below ERM, reaching $-0.010$ at $2000$. The oracle-to-ERM headroom
persists across the same range, so more data does not convert into recovered
performance. Calibration behaves in the opposite way (Fig.~\ref{fig:nsweep}b):
under label-prior shift the recovery achieved by prior adjustment grows with the
amount of unlabeled target data, from $0.10$ at $250$ patients to $0.18$ at
$2000$, because the target prevalence is estimated from those unlabeled
covariates. The two axes thus have different sample-size signatures.

\textbf{Nonlinearity.}
The negative result is not an artifact of a linear observation model. Replacing
the linear observation map with a saturating variant ($\alpha=0.5$) adds a
nonlinear channel that flexible neural methods could in principle exploit.
Across the same sample-size sweep, adaptation still recovers at most $0.005$
AUROC over ERM.

\textbf{Limitations.}
The controlled results rest on synthetic data whose ground truth we specify.
This is what makes the two references and a falsifiable ceiling possible, but it
also means the generator is linear and low-dimensional, and that each scenario
perturbs a single mechanism whereas real shift combines several. This simplicity
works against the methods rather than for them: the invariant subspace exists by
construction and is still not recovered.

The real cohorts are a cautious counterpart rather than independent
confirmation. The I3LUNG transfers share source centers across pairs and the
registry contributes only two centers, so the pooled intervals are mildly
optimistic; we report a single horizon for the registry for this reason. On real
data the shift mechanisms are largely latent: only covariate and label-prior
shift are measurable from unlabeled data, and the remaining mechanisms must be
interpreted through the benchmark rather than identified directly.

Finally, we study a binary fixed-horizon outcome at a single horizon $\tau$
rather than the full survival function, and evaluate a representative but not
exhaustive set of methods. Extending the benchmark to censored time-to-event
losses and to generative adaptation methods is left to future work.
% ==== END Ablation ====================================================


% ==== REPLACE: \section{Discussion} (intero) ==========================
\section{Discussion}

The results describe what unsupervised transfer can and cannot achieve in a
small-sample, multi-center regime, and attribute each outcome to a generative
cause. Discrimination is not recovered by any method we evaluated, even where
the observed-data oracle shows the signal is still present in the target
covariates. Because the shortfall persists as the sample grows, it reflects an
identifiability limit rather than a shortage of data: without target labels,
the source-trained decision rule cannot be told apart from the alternatives the
unlabeled target covariates are equally consistent with. Calibration behaves
differently. It is recoverable without labels when prevalence is the only
quantity that changes, the recovery sharpens as more unlabeled target data
arrives, and it degrades when the label-shift assumption is violated.

This has a direct consequence for how a transfer should be approached in a
multi-center clinical project. The two shift mechanisms measurable without
target labels---a change in the covariate distribution and a change in event
prevalence---are also the two that behave most differently, which makes them a
usable decision rule. Before deploying a model at a new center, measure both,
since both are computable from unlabeled data. If prevalence has moved,
recalibrate: that is the one regime in which an unsupervised correction
demonstrably restores the probabilities a threshold-based decision depends on.
If discrimination has dropped, our results indicate that no amount of
unsupervised alignment will return it, and the effort is better spent upstream
of the model---harmonizing how predictors are measured across centers, or
acquiring labels at the target site. The contribution we intend is diagnostic
rather than prescriptive: a reported transfer failure conflates a transfer that
is intrinsically unsolvable from unlabeled data with a method that
underperforms on a solvable one, and separating the two is what tells a clinical
team whether to change the method or to change the data.

A further extension concerns modality. We deliberately study tabular clinical
data: it is the most widely available modality across centers, the one on which
most deployed risk models are built, and the setting in which invariance
assumptions can be stated and tested most cleanly. Oncology is nonetheless
moving toward decisions informed by imaging, digital pathology, and molecular
profiling alongside the clinical record, and multimodal foundation models are
increasingly proposed to combine
them~\cite{wang2024pathologyfm,xiang2025visionlanguage,ding2025wsifm}.
Multi-center shift is not removed by this move and may be sharpened by it, since
scanners, staining protocols, and sequencing pipelines differ across sites in
ways that clinical variables do not. Whether a pretrained representation
transports the invariant signal that alignment fails to recover here is an open
question, and the diagnostic framing of this paper---a known oracle with one
mechanism perturbed at a time---transfers to that setting unchanged, since it
constrains the evaluation design rather than the model class.
% ==== END Discussion ==================================================


% ==== INSERT: in §Results, real data, dopo la frase sul best-per-target
% Rende esplicito il risultato significativo ma non rilevante in LOCO.

Two summaries do reach nominal significance in the
domain-generalization protocol---IW-ERM at $+0.0062$ ($p=0.004$) and the
best-per-fold selection at $+0.0072$ ($p=0.033$)---but both are below a third of
the pre-specified $0.02$ margin, and the best-per-fold figure is an optimistic
selection made on target performance. We therefore read them as statistically
detectable but not clinically relevant.
% ==== END ============================================================


% ==== BIBTEX: aggiungere a references.bib ============================
% Verificate 2026-08-19 via Undermind (DOI controllati).

@article{wang2024pathologyfm,
  author  = {Wang, Xiyue and others},
  title   = {A pathology foundation model for cancer diagnosis and prognosis
             prediction},
  journal = {Nature},
  volume  = {634},
  pages   = {970--978},
  year    = {2024},
  doi     = {10.1038/s41586-024-07894-z}
}

@article{xiang2025visionlanguage,
  author  = {Xiang, Jinxi and others},
  title   = {A vision--language foundation model for precision oncology},
  journal = {Nature},
  volume  = {638},
  pages   = {769--778},
  year    = {2025},
  doi     = {10.1038/s41586-024-08378-w}
}

@article{ding2025wsifm,
  author  = {Ding, Tong and others},
  title   = {A multimodal whole-slide foundation model for pathology},
  journal = {Nature Medicine},
  volume  = {31},
  pages   = {3749--3761},
  year    = {2025},
  doi     = {10.1038/s41591-025-03982-3}
}
% ==== END bibtex =====================================================
